\documentclass[12pt]{article}
\usepackage[utf8]{inputenc}
\usepackage[T1]{fontenc}
\usepackage[margin=1in]{geometry}
\usepackage{setspace}
\usepackage{xcolor}
\usepackage[colorlinks=true, urlcolor=blue!60!black, linkcolor=blue!60!black]{hyperref}
\setlength{\parindent}{2em}
\setlength{\parskip}{1.25em}
\renewcommand{\baselinestretch}{1.0}
\begin{document}
\begin{center}
{\LARGE \textit{United States v. Pink}, 315 U.S. 203 (1942)\par}
\vspace{0.8em}
{\large Case Brief --- Four-Page Edition\par}
\end{center}
\vspace{1em}

\section{Memory Jogger}

The Litvinov Assignment, an executive agreement the President made without the Senate when recognizing the Soviet Union, gave the United States title to a nationalized Russian insurer's New York assets; the Court held 5--2 that the agreement overrode New York's contrary law and policy just as a treaty would (\href{https://supreme.justia.com/cases/federal/us/315/203/}{315 U.S.\ 203}).

\section{Detailed Case Facts}

The First Russian Insurance Company, organized under the laws of the Russian Empire, opened a New York branch in 1907 and deposited assets with New York authorities to secure local policyholders and creditors. After the Revolution, Soviet decrees of 1918 and 1919 nationalized all Russian insurers, purporting to extinguish their debts and confiscate their assets wherever located.

In 1925 New York placed the branch in liquidation under its Superintendent of Insurance, eventually Louis H.\ Pink. After paying all domestic creditors, the Superintendent held a surplus of more than \$1 million, and in 1931 the New York Court of Appeals directed him to pay foreign creditors whose claims arose abroad and remit the rest to the company's surviving directors. Before distribution was complete, President Roosevelt recognized the Soviet Union on November 16, 1933. As part of recognition, and without Senate ratification, he concluded the Litvinov Assignment, under which the Soviet government assigned to the United States all amounts due it from American nationals, creating a fund for American citizens' confiscation claims.

\section{Procedural History}

In 1934 the United States sued Pink and designated foreign creditors in New York Supreme Court to recover the surplus. Relying on \textit{Moscow Fire Insurance Co.\ v.\ Bank of New York \& Trust Co.}, which the Supreme Court had affirmed by an equally divided vote, the New York courts dismissed, holding that the Soviet decrees had no extraterritorial effect and that New York's policy against confiscation governed. The Court granted certiorari, heard argument on December 15, 1941, and reversed on February 2, 1942.

\section{Judicial Votes}

Reversed and remanded, \textbf{5--2}.
\begin{itemize}
\item \textbf{Majority:} Justice Douglas, for Justices Black, Frankfurter, Murphy, and Byrnes.
\item \textbf{Concurring:} Justice Frankfurter.
\item \textbf{Dissenting:} Chief Justice Stone, joined by Justice Roberts.
\item \textbf{Not participating:} Justices Reed and Jackson.
\end{itemize}

\section{Holding}

Extending \textit{United States v.\ Belmont} (1937), the Court held that the President's power to recognize foreign governments includes the implied power to settle outstanding claims by a sole executive agreement, and that such an agreement carries the same supremacy over inconsistent state law and policy as an Article II treaty. The 1918--1919 decrees reached the New York surplus, title passed to the United States, and neither New York's marshaling decree nor the Fifth Amendment barred the United States from taking priority over foreign creditors whose claims arose abroad.

\section{Analysis of Opinions}

\textbf{Majority (Douglas).} The 4--4 affirmance in \textit{Moscow Fire} had no precedential force, and because the federal right depended on the reach of the Soviet decrees, the Court could determine Russian law itself; relying on an official Soviet declaration, it found the decrees meant to reach foreign assets. Although aliens enjoy Fifth Amendment protection, due process does not bar the United States from preferring itself and its nationals over foreign creditors whose claims did not arise from the New York branch's business. Recognition and claims settlement were parts of one diplomatic bargain, and under \textit{Curtiss-Wright} removing obstacles to recognition was a ``modest implied power'' of the President as ``sole organ'' in foreign relations. Because foreign-affairs power belongs exclusively to the national government, sole executive agreements share treaties' supremacy over state law, and letting New York refuse enforcement would revive the very friction the President had resolved.

\textbf{Frankfurter (concurring).} Domestic conflict-of-laws concepts such as ``situs,'' ``jurisdiction,'' and ``comity'' should not govern international controversies; once the political branch accepted the Assignment as the instrument of settlement, state courts could not frustrate it through local public policy.

\textbf{Stone (dissenting).} Recognition validates a regime's acts only within its own territory and does not retroactively override New York property law or transfer New York assets already marshaled for creditors. Even if the President could displace state law expressly, the Assignment transferred only whatever rights the Soviet Union actually had and contained no command overriding state priority rules or extinguishing third-party claims.

\textbf{Later developments.} \textit{Dames \& Moore v.\ Regan} (1981) upheld the Algiers Accords, relying on congressional acquiescence alongside \textit{Belmont} and \textit{Pink}. \textit{American Insurance Ass'n v.\ Garamendi} (2003) held that Holocaust-era insurance agreements preempted a California statute. \textit{Medell\'in v.\ Texas} (2008) confined these cases to a ``narrow set of circumstances'' involving claims settlement with congressional acquiescence.

\section{Examples: Future Applications}

\begin{enumerate}
\item \textbf{Same-Side (Preempts): Nation Z's Airline Assets.} An agreement incident to recognition assigning a nationalized airline's U.S.\ assets to the United States for American claimants is classic claims settlement and preempts state distribution rules.
\item \textbf{Same-Side (Preempts): Releasing Attachments.} An exchange of letters sending American claims to a commission and releasing state-court attachments of Nation R's deposits binds Illinois courts.
\item \textbf{Opposite-Side (No Preemption): Limitations Waivers.} Requiring states to waive tort limitations periods for Nation Z's diplomats lacks the pedigree of claims settlement and falls on the \textit{Medell\'in} side.
\item \textbf{Opposite-Side (Likely Ineffective): Overriding Congress.} An agreement exempting Nation W's firms from statutory sanctions puts the President at his lowest ebb under \textit{Youngstown}; \textit{Pink} addressed only state law.
\item \textbf{Fence-Sitter: Extinguishing American Liens.} A lump-sum settlement transferring frozen assets that American creditors had attached raises a takings question \textit{Pink} did not face and \textit{Dames \& Moore} left open.
\end{enumerate}

\section{Critique}

\textbf{Textual puzzle.} The Supremacy Clause lists the ``Constitution,'' ``Laws of the United States,'' and ``Treaties,'' not sole executive agreements, so \textit{Pink}'s equation of the two invites textualist objection.

\textbf{Federalism and separation of powers.} Stone's dissent presses a further concern: if the President can bypass the Senate's two-thirds requirement and divest third parties' state-law property rights through an exchange of letters, little would stop the Executive from using foreign-policy agreements to rewrite domestic private law. \textit{Medell\'in}'s narrowing reflects that worry.

\section{Key Quotations}

\begin{quote}
``Power to remove such obstacles to full recognition as settlement of claims of our nationals ... certainly is a modest implied power of the President who is the `sole organ of the federal government in the field of international relations.'\,''
\end{quote}
\noindent --- Justice Douglas, for the Court.

\begin{quote}
``Enforcement of such state policies would indeed tend to restore some of the precise impediments to friendly relations which the President intended to remove on inauguration of the policy of recognition of the Soviet Government.''
\end{quote}
\noindent --- Justice Douglas, for the Court.

\begin{quote}
``Legal ideas like other organisms cannot survive severance from their congenial environment. Concepts like `situs' and `jurisdiction' and `comity' summarize views evolved by the judicial process, in the absence of controlling legislation, for the settlement of domestic issues. To utilize such concepts for the solution of controversies international in nature, even though they are presented to the courts in the form of a private litigation, is to invoke a narrow and inadmissible frame of reference.''
\end{quote}
\noindent --- Justice Frankfurter, concurring.

\end{document}
